\section{Subexponential skeleton cost and the large repeat tail}\label{sec:uppercost}
We now show that the extra skeleton information in \eqref{eq:finiteupper} does not change the linear logarithmic constant. Let $L=\log n$ and
\[
 R=\left\lceil\frac{4n\log L}{L}\right\rceil.
\]
Eventually $R<n/2$. For $1\le r\le R$, the binomial coefficients $\binom nh$ are increasing for $h\le r$. Hence
\[
 \sum_{h=0}^r\binom{n-1}{h}\le(r+1)\binom nr.
\]
Using $\binom ab\le(ea/b)^b$ in the three factors of \eqref{eq:skeleton} yields
\begin{equation}\label{eq:smallS}
 S(n,r)\le(r+1)(2en/r)^{3r},\qquad
 \max_{0\le r\le R}\log S(n,r)=O\!\left(\frac{n(\log L)^2}{L}\right).
\end{equation}
For the second estimate, $r\log(2en/r)$ is increasing for $r\le n$, and at $r=R$ the factor $\log(2en/r)$ is $O(\log L)$. The case $r=0$ has zero skeleton cost. Consequently the small-$r$ part is bounded by
\begin{equation}\label{eq:smallupper}
 \sum_{r=0}^R S(n,r)T(n+1,r+1)
 \le\exp(O(E(n)))b_{n+1}.
\end{equation}

A uniform estimate for the remaining range is essential: it is not valid to assume the encoding's $r$ lies near a triangular saddle. From \eqref{eq:skeleton}, for every $r<n$,
\[
 S(n,r)\le2^n\,2^{2n}\,2^n=16^n.
\]
For $n\ge2$, the top argument in \eqref{eq:vandermonde} is at most $n^2$. Write $K=n-r\ge1$. Then
\begin{align*}
 \log T(n+1,r+1)
 &\le K\log(en^2/K)\\
 &=(n-r)L+K+K\log(n/K)\\
 &\le(n-r)L+(1+1/e)n.
\end{align*}
The last step uses $0<t\le1\Rightarrow t\log(1/t)\le1/e$. For $r>R$, $rL\ge4n\log L$, and the sum of at most $n$ such terms is therefore at most
\begin{equation}\label{eq:tail}
 \exp\bigl(nL-4n\log L+O(n)\bigr).
\end{equation}
Lemma~\ref{lem:triangular} already gives a single triangular term of logarithm $nL-2n\log L+O(n)$, so this tail is smaller by more than any fixed linear exponential factor eventually.

Combining \eqref{eq:smallupper}, \eqref{eq:tail}, and Lemma~\ref{lem:triangular}, and observing that replacing $n$ by $n+1$ changes the displayed main term by $O(L)$, gives
\begin{equation}\label{eq:upperfinal}
 \log u_n\le n(L-2\log L+C_+)+O(E(n))
 \qquad(u=a,d).
\end{equation}
The same upper bound follows for $e,c$ by containment in $a,d$. It remains to produce enough ordinary words, with losses explicitly smaller than $\exp(\varepsilon n)$ for every fixed $\varepsilon>0$.
